% Generated draft: requires mathematical and source review.
\documentclass[11pt]{article}
\usepackage[margin=25mm]{geometry}
\usepackage{fontspec}
\setmainfont{Latin Modern Roman}
\usepackage{amsmath,amssymb,mathtools,unicode-math}
\setmathfont{Latin Modern Math}
\usepackage{enumitem,needspace,xurl,hyperref}
\setlength{\parindent}{0pt}
\setlength{\parskip}{5pt}
\setlength{\emergencystretch}{3em}
\newcommand{\N}{\mathbb{N}}
\newcommand{\Z}{\mathbb{Z}}
\newcommand{\Q}{\mathbb{Q}}
\newcommand{\R}{\mathbb{R}}
\newcommand{\C}{\mathbb{C}}
\newcommand{\F}{\mathbb{F}}
\newcommand{\A}{\mathbb{A}}
\newcommand{\PP}{\mathbb{P}}
\newcommand{\E}{\mathbb{E}}
\newcommand{\eps}{\varepsilon}
\newcommand{\dd}{\,\mathrm{d}}
\begin{document}
\section*{UnsolvedMath ID: 700604}
\textbf{Source problem code:} catalogue entry
\subsection{Definitions and mathematical statement}
Let \(\Omega\) be an infinite set and let \(G\leq\operatorname{Sym}(\Omega)\) be a permutation group. The action is \emph{2-transitive} if, for any \(\alpha,\beta,\gamma,\delta\in\Omega\) with \(\alpha\neq\beta\) and \(\gamma\neq\delta\), there exists \(g\in G\) such that \(g(\alpha)=\gamma\) and \(g(\beta)=\delta\). For \(\alpha\in\Omega\), write
\[
G_\alpha=\{g\in G:g(\alpha)=\alpha\}.
\]
A group is \emph{locally soluble} if every finitely generated subgroup is soluble. Here a group \(H\) is soluble if its derived series, defined by \(H^{(0)}=H\) and \(H^{(j+1)}=[H^{(j)},H^{(j)}]\), reaches the identity subgroup after finitely many steps.

\textbf{Question.} Suppose that \(G\) is 2-transitive and that \(G_\alpha\) is locally soluble for some \(\alpha\in\Omega\). Must there exist a field \(K\), an integer \(n\geq1\), and a homomorphism
\[
\rho:G\longrightarrow\operatorname{GL}_n(K)
\]
such that \(\rho(g)\neq I_n\) for at least one \(g\in G\), and the only \(\rho(G)\)-invariant \(K\)-linear subspaces of \(K^n\) are \(\{0\}\) and \(K^n\)?

Point stabilizers in a transitive action are conjugate, so the local solubility hypothesis then holds at every point. We use the usual faithful meaning of permutation group; under this convention, an infinite permutation group cannot have finite degree. The field and dimension may depend on \(G\). Non-triviality does not require faithfulness of \(\rho\), and irreducibility is over \(K\), without an additional requirement of absolute irreducibility.
\subsection{Short English statement}
Does every infinite 2-transitive permutation group with locally soluble point stabilizers admit a non-trivial irreducible representation of finite positive dimension over some field? This is the question attributed to A. N. Fomin in Kourovka Notebook Problem 9.71.
\subsection{Sources}
\begin{itemize}
\item Supplied local UnsolvedMath registry record: permanent ID \(700604\), exact source problem code \texttt{catalogue entry}, title \emph{Kourovka Notebook Problem 9.71}; upstream revision \texttt{d371ee917ac9647fbb06cd6fc529deec6f3dda82}. The record was created and updated on 2026-10-16 and labels the problem as open; this is supplied registry information, not an independently verified status assessment.
\item A. N. Fomin, Problem 9.71, \emph{The Kourovka Notebook (Unsolved Problems in Group Theory)}, 9th issue (1984), as attributed in the supplied record.
\item The supplied background also identifies the 21st issue edition (2026), \texttt{arXiv:1401.0300v48}. The collection description names E. I. Khukhro and V. D. Mazurov as editors. This bibliographic attribution is supplied source information and has not been independently checked.
\end{itemize}
\end{document}
